

\section{Benutzungshinweise, Abkürzungsverzeichnis und Nomenklatur}

% \Text
\paragraph{Maßnahmenkatalog}
{Ergänzend zu diesem Lehrbuch existiert ein eigener
Maßnahmenkatalog, in welchem die jeweiligen konkreten
Maßnahmen zu finden sind. Referenzen mit dem Symbol
{\SignMassnahme} beziehen sich auf diesen Katalog. Sofern
nicht anders angegeben, beziehen sich die Angaben auf den 
Maßnahmenkatalog in der \EE{Version \VarMkatalogVersion}.

}
{}
{}
{}
{}

\Text{Externe Teile}
{Aus Kostengründen sind manchmal nicht
alle Teile Bestandteil eines Drucks. Dies betrifft
insbesonders Teile des Appendix wie z.\,B. das
Literaturverzeichnis oder den Kommentarteil. Diese Teile
sind online auf der {\AASS}-Projektseite im
Download-Bereich verfügbar
(\url{http://www.aass.at/cms/download/aass/}).

}

\Text{Im Text}{%
\E{Betonte} Textteile sind kursiv, \EE{Besonders stark
betonte} sind fett gedruckt. \T[noindex]{Fachbegriffe}, welche an
dieser Stelle erklärt werden, sind ebenso fett gedruckt.
\Z{Zusatzinformationen}, welche der Vollständigkeit halber
niedergeschrieben sind, werden grau gedruckt: Sie dienen der
zusätzlichen Informationen und stellen keinen inhaltlichen
Schwerpunkt dar. 

Ähnliches gilt für Fußnoten: Normale
Fußnoten\footnote{Ich bin eine normale Fußnote.} beinhalten
Hintergrundinformationen, die nicht uninteressant sind, aber
den normalen Fließtext  unnötig aufblähen würden.
Zusatzinfo-Fußnoten\ZFootnote{Ich bin eine
Zusatz-Info-Fußnote.} werden grau gedruckt und beinhalten
Hintergrundinformationen für besonders interessierte Leser.

An manchen Stellen im Text befindet sich eine
Kennzeichenungen wie z.\,B.
\Commentary[Benutzungshinweise]{Ich bin ein
Beispielkommentar aus dem Einleitungsteil.}. Dies ist ein
Hinweis auf ein Kommentar im Anhang (Appendix).

}

% 
% \Text{Index}{%
% \begin{description}
%   \item[Zusammengesetzte Begriffe] Eigenschaftswort
%   Hauptwort \Sign{→} Hauptwort, Eigenschaftswort; \enquote{Begriffe,
%   zusammengesetzte}
%   \item[Umlaute] Ä, Ö, Ü \Sign{→} Ae, Oe, Ue
%   \item[Lateinische Begriffe] Lateinische Hauptwörter werden im Index groß geschrieben.
% \end{description}
% }{}{}{}{}

\Text{Kompetenzlevel}
{}

% \TableWide
% 	{Titel}%1
% 	{Beschreibung}%2
% 	{Label}%3
% 	{Quelle}%4
% 	{Lizenz}%5
% 	{Tabelle}%6
% 	{Float}%8
\TableWide
{Kompetenzlevel}%1
{}%2
{}%3
{}%4
{}%5
{%
\begin{tabularx}{\linewidth}{cclcX}
\addlinespace\toprule\addlinespace
\noindent\TabHeading{Level}
	& \noindent\TabHeading{Level (Text)}
					& \noindent\TabHeading{Titel}
													& \noindent\TabHeading{ALS-Äquivalent}
															& \noindent\TabHeading{Beschreibung}
\\\addlinespace\midrule\addlinespace
0 	& \TabSub{A}	& Laien, Ersthelfer				& 		& Laienhelfer
\\\addlinespace
1	& \TabSub{B} 	& First Responder				& AED	& First Responder (ohne Ausbildung nach SanG)
\\\addlinespace
2	& \TabSub{C} 	& Basic							& BLS	& \Lang{.at} Rettungssanitäter
\\\addlinespace
3	& \TabSub{D} 	& Intermediate					& ILS	& \Lang{.at} Notfallsanitäter mit oder ohne Notfallkompetenz NKA; \newline
															  \Lang{.de} Rettungssanitäter
\\\addlinespace
4	& \TabSub{E} 	& Advanced						& ALS   & Ärzte, diplomiertes Pflegepersonal; \newline
															  \Lang{.at} Notfallsanitäter mit Notfallkompetenz NKV oder NKI; \newline
															  \Lang{.de} Notfallsanitäter, Rettungsassistenten \newline
\\\addlinespace
5	& \TabSub{F} 	& Expert						&		& Experten, Ärzte, Fachärzte, diplomiertes Pflegepersonal mit Zusatzausbildung
\\\addlinespace\bottomrule\addlinespace
\end{tabularx}

}%6
{H}%8

% \TableWide
% 	{Titel}%1
% 	{Beschreibung}%2
% 	{Label}%3
% 	{Quelle}%4
% 	{Lizenz}%5
% 	{Tabelle}%6
% 	{Float}%8
\Table
{Maßnahmentypen}%1
{}%2
{}%3
{}%4
{}%5
{%
\begin{tabularx}{\linewidth}{clX}
\addlinespace\toprule\addlinespace
\noindent\TabHeading{ID}
	& \noindent\TabHeading{Art}
										& \noindent\TabHeading{Beschreibung}
\\\addlinespace\midrule\addlinespace
ALL & \TabSub{Allgemeine Maßnahmen}		& Allgemeine Maßnahmen beziehen sich auf eine Gruppe von Patienten mit einer Übergruppe bestimmter, verwandter Erkrankungen, wobei eine weiterführende Unterteilung mittels spezieller Maßnahmen existiert.			
\\\addlinespace
STD	& \TabSub{Standardmaßnahmen} 		& Standardmaßnahmen treffen unter bestimmten Bedingungen auf eine große Gruppe unterschiedlicher Patienten mit unterschiedlichen Diagnosen zu.
\\\addlinespace
SPZ	& \TabSub{Spezielle Maßnahmen}	 	& Spezielle Maßnahmen treffen auf bestimmte Patienten mit einer bestimmten Diagnose oder einem bestimmten Zustand zu.						
\\\addlinespace\bottomrule\addlinespace
\end{tabularx}

}%6
{H}%8

\begin{WidePar}
\begin{multicols}{3}
{\raggedright\ColumnFontNarrow%
\minisec{Symbole und Einheiten}

\begin{labeling}{mmm}
  \item[\Female] Frau, weiblich
  \item[\Male] Mann, männlich
  \item[\Min] Minute 
  \item[\Sek] Sekunde  
  
  \item[\Ca{!}] Cave
  \item[\SymbolDefinition] Definition
  \item[\SignFigure] Abbildung
  \item[\libertineGlyph{creativecommons}] Creative Commons
  \item[\SignFazit] Leitsatz
  \item[\SignMassnahme] Maßnahmen, Referenzen beziehen
  sich auf den Maßnahmenkatalog
  \item[\SignTable] Tabelle, Tafel
  \item[\dsliterary] Literaturverweise
\end{labeling}

\minisec{Abkürzungen}
\begin{description}
% \setlength{\itemindent}{0pt}
% \setlength{\listparindent}{0pt}
% \setlength{\leftmargin}{0pt}
% \setlength{\labelwidth}{0pt}
\setlength{\parskip}{0pt}%
  \item[A.] \Lang{lat.} Arteria, 
  Arterie\index{A.|see{Arteria}}
  \item[AASS]  Arbeits- und Ausbildungsstandards für den
  Sanitätdienst.
  \item[ABGB]   Allgemeines Bürgerliches Gesetzbuch
  \item[Abk.]   Abkürzung
  \item[Abs.]   Absatz
  \item[AF]     Atemfrequenz
  \item[AIDS]   Acquired Immune Deficiency Syndrome
  \item[allg.]  allgemein
  \item[AMV]    Atemminutenvolumen
  \item[ARGE]   Arbeitsgemeinschaft
  \item[ASVG]   Arbeits- und Sozialversicherungsgesetz
  \item[AUVA]   Allgemeine Unfallversicherungsanstalt
  \item[AO]     Abgabeort
  \item[AR]     Arbeitsrecht
  \item[AZV]    Atemzugsvolumen
  \item[B-VG]   Bundesverfassungsgesetz
  \item[BKA]    Bundeskanzleramt, Bundeskriminalamt
  \item[BM]     Bundesministerium
  \item[BMI]    Bundesministerium für Inneres
  \item[BO]     Berufungsort
  \item[BGBl]   Bundesgesetzblatt
  \item[BP]		\Lang{engl.} blood pressure: Blutdruck
  \item[Cave]   \Lang{lat.} \Speech{\enquote{Hüte Dich vor
  \dots}}. Auch: \enquote{Achtung, gib Acht!}
  \item[CICR]	\Lang{franz.} Comité international de la
  Croix-Rouge: Internationales Komittee des
  Roten Kreuz
  \item[CPR]    Cardio-Pulmonale Reanimation
  \item[DNHG]   Dienstnehmerhaftpflichtgesetz
  \item[EGC]    Europäische Grundrechtecharta
  \item[EMRK]   Europäische Menschenrechtskonvention
  \item[engl.]  englisch
  \item[ETI]    endotracheale Intubation
  \item[evtl.]  eventuell
  \item[EU]     Europäische Union
  \item[fETI]   frühe endotracheale Intubation
  \item[franz.] französisch
  \item[FW]     Feuerwehr
  \item[griech.] griechisch.
  \item[gt.]    \Lang{lat., Abkz.} gutta: der Tropfen; \Lang{Pl.} gtt.
  (guttae): die Tropfen. (Nicht genormte) Maßeinheit.
  \item[Gz]     Geschäftszahl
  \item[HB]     Heilbehandlung
  \item[HDM]    Herzdruckmassage
  \item[HF]     Herzfrequenz
  \item[HIV]    Human Immundeficiency Virus
  \item[i.\,d.\,R.] in der Regel
  \item[i.\,V.\,m.] in Verbindung mit
  \item[KA]     Krankenanstalt
  \item[KAKuG]  Krankenanstalten- und Kuranstaltengesetz.
  \item[KARPAT] Pathologisch-Bakteriologischen Institut der Krankananstalt
  Rudolfstiftung der Stadt Wien
  \item[KFG]    Kraftfahrgesetz
  \item[KG]     Körpergewicht
  \item[KH]     Krankenhaus
  \item[KHD]    Katastrophen- und Großschadenshilfsdienst
  \item[KIT]    Kriseninterventionsteam
  \item[KTW]    Krankentransportwagen
  \item[lat.]   latein, lateinisch
  \item[LT]		Larynxtubus
  \item[MA]     Magistratsabteilung
  \item[MCI]    Myocardinfarkt, Herzinfarkt
  \item[mm\,Hg] Millimeter
  Quecksilbersäule
  \item[MONA]   Behandlungsschema für Patienten mit
  Herzinfarkt, bestehend aus Morphin, \ce{O2}-Gabe, Nitrat-Gabe (z.\,B.
  Nitrolingual) und Gabe von Acetylsalicylsäure (Aspirin{\texttrademark},
  Aspisol{\texttrademark})
  \item[MPG]    Medizinproduktegesetz
  \item[N.]
  \begin{inparaenum}
    \item \Lang{lat.} Nervus: Nerv\index{N.|see{Nervus}}
    \item \Lang{lat.} Neoplasma: Neubildung\index{N.|see{Neoplasma}}
  \end{inparaenum}
  \item[NA]     Notarzt
  \item[\ce{NaCl}] Natrium-Chlorid
  \item[NEF]    Notarzteinsatzfahrzeug
  \item[neurol.] neurologisch; das Nervensystem betreffend.
  \item[NFS]    Notfallsanitäter
  \item[NKA]    Notfallsanitäter mit allgemeiner Notkompetenz Arzneimittel.
  \item[NKI]    Notfallsanitäter mit besonderer Notkompetenz Intubation und Beatmung
  \item[NKV]    Notfallsanitäter mit allgemeiner Notkompetenz venöser Zugang und Infusion
  \item[NKx]    Notfallsanitäter mit Notfallkompetenzen, ohne nähere
  Spezifizierung welche Notfallkompetenzen gemeint sind
  \item[OGH]    Oberster Gerichtshof
  \item[OLG]    Oberlandesgericht
  \item[OSR]    Oberster Sanitätsrat
  \item[PatVG]  Patientenverfügungsgesetz
  \item[PD]     Public domain
  \item[PF]		Pulsfrequenz
  \item[PI]     Polizeiinspektion
  \item[Pl.]    Plural
  \item[PLS]    Patientenleitsystem
  \item[PLT]    Patientenleittasche
  \item[RM]     Rettungsmittel
  \item[RQW]    Rissquetschwunde
  \item[RR]     Blutdruck nach Riva Rocci, oft allgemein für Blutdruck
  \item[RS]     Rettungssanitäter
  \item[RTW]    Rettungstransportwagen
  \item[Rz.]    Randzahl
  \item[S.]     Seite
  \item[SanAVO] Sani\-tät\-er-Aus\-bildungs\-ver\-ord\-nung.
  \item[SanG]   Sanitätergesetz.
  \item[scherzh.] scherzhaft
  \item[SHT]    Schädel-Hirn-Trauma.
  \item[SMG]    Suchtmittelgesetz
  \item[SMZ]	Sozialmedizinisches Zentrum
  \item[spez.]  speziell
  \item[StGB]   Strafgesetzbuch
  \item[StGG]   Staatsgrundgesetz
  \item[StPO]   Strafprozessordnung
  \item[StrR]   Strafrecht
  \item[StVO]   Straßenverkehrsordnung
  \item[SV]     Sozialversicherung
  \item[Syn.]   Synonym
  \item[UbG]    Unterbringungsgesetz
  \item[ugs.]   umgangssprachlich
  \item[USAF]   United States Air Force
  \item[USMC]   United States Marine Corps
  \item[urspr.] ursprünglich
  \item[VfGH]   Verfassungsgerichtshof
  \item[VfR]    Verfassungsrecht
  \item[VwGH]   Verwaltungsgerichtshof
  \item[VwR]    Verwaltungsrecht
  \item[WmCo]   Wikimedia Commons
  \item[Z.]     Zahl, Ziffer.
  \item[ZDG]    Zivildienstgesetz
  \item[ZDL]    Zivildienstleistender
  \item[ZR]     Zivilrecht
\end{description}
}
\end{multicols}
\end{WidePar}